\documentclass[11pt]{article}

\usepackage[margin=1in]{geometry}
\usepackage{amsmath,amssymb,mathtools}
\usepackage[inline,shortlabels]{enumitem}
\usepackage{booktabs,array}
\usepackage{xcolor}
\usepackage{microtype}
\usepackage[colorlinks=true,linkcolor=blue!60!black,
  urlcolor=blue!60!black]{hyperref}
\usepackage[nameinlink,capitalize,noabbrev]{cleveref}

\definecolor{courseblue}{HTML}{17365D}
\newcommand{\E}{\mathbb E}
\newcommand{\I}{\mathbb I}
\renewcommand{\P}{\mathbb P}
\newcommand{\Law}{\mathcal L}
\newcommand{\R}{\mathbb R}

\hypersetup{
  pdftitle={CMPUT 654 probability notation conventions},
  pdfauthor={Csaba Szepesvari}
}

\setlength{\parindent}{0pt}
\setlength{\parskip}{0.55em}

\begin{document}

\begin{center}
\fcolorbox{courseblue}{white}{%
\begin{minipage}{0.92\textwidth}
\vspace{1mm}
\textbf{CMPUT 654: Mathematical Foundations of Modern AI Systems}
\hfill Fall 2026

\vspace{3mm}
\begin{center}
{\Large\bfseries Probability notation conventions}

\vspace{1mm}
Course reference for lecture notes and solutions
\end{center}

\vspace{2mm}
\textit{Last revised: September 27, 2026}
\vspace{1mm}
\end{minipage}}
\end{center}

\section{Purpose}

Probability notation is meaningful only after the types of its objects are
clear. These conventions distinguish an underlying probability measure, a
random element, its law, a probability mass function, and a predictive
kernel. They also specify the shorthand accepted in this course.

Standard probability notation lets us state the experiment once and then
calculate without repeatedly annotating every probability and expectation.
Inconsistent shorthand common in machine-learning writing can use the same
expression for a probability, a law, and a density, or leave the source of
randomness implicit. Such notation becomes ambiguous as soon as a calculation
contains data, initialization, and algorithmic randomness. The distinctions
below prevent this ambiguity.

The conventions apply to lecture notes, exercises, solutions, and homework.
They are defaults for mathematical writing in the course; a different
convention must be introduced explicitly and used consistently.

\section{The main objects have different types}

An experiment has a set \(\Omega\) of possible outcomes. An \emph{event}
\(E\) is a subset of \(\Omega\), such as the set of outcomes on which the
prediction is correct. The probability \(\P(E)\) is a number assigned to this
event.

A \emph{random variable}, or more generally a \emph{random element}, is a map
\[
 Z:\Omega\longrightarrow\mathcal Z.
\]
It assigns a value \(Z(\omega)\in\mathcal Z\) to each possible outcome
\(\omega\in\Omega\). The training sample and the learned predictor are random
elements when they depend on the outcome of the experiment. A random element
\(Z\) has a \emph{law}, written \(\Law(Z)\), which assigns probabilities to
sets of possible values in \(\mathcal Z\).\footnote{There are technical
restrictions on the permitted subsets and maps. If \(\Omega\) is finite or
countable, every subset may be declared an event, and a probability law is
specified by assigning masses to the individual outcomes. The issue arises
for uncountable spaces such as \(\mathbb R\). For example, length on
\([0,1]\) cannot be extended to every subset while preserving countable
additivity and translation invariance. One therefore specifies a
sigma-algebra \(\mathcal G\) of permitted events in \(\Omega\), a
sigma-algebra \(\mathcal H\) of permitted subsets of \(\mathcal Z\), and
requires \(Z\) to be measurable: \(\{Z\in A\}\in\mathcal G\) for every
\(A\in\mathcal H\). These restrictions do not change the basic types used
here: events are subsets of \(\Omega\), random elements are maps from
\(\Omega\), and laws assign probabilities to sets of values.}

Probability measures, laws, kernels, and expectations are built from these
basic types. The following table records their input and output types.

\begin{center}
\renewcommand{\arraystretch}{1.25}
\begin{tabular}{>{\(}l<{\)}p{0.36\textwidth}p{0.34\textwidth}}
\toprule
\text{Notation} & \text{Argument} & \text{Value}\\
\midrule
\P & an event & a number in \([0,1]\)\\
\Law & a random element or tuple of random elements & a probability law\\
P,Q & as declared; commonly a set supplied to a law, or a value and input supplied to a kernel & a number in \([0,1]\)\\
\E,\E_P & an integrable random variable, or an integrable function when the law \(P\) is named & a real number\\
p,q & as declared; commonly a mass function, density, probability vector, or kernel & a number or probability vector\\
\bottomrule
\end{tabular}
\end{center}

The fundamental relation between the ambient probability measure and a law is
\begin{equation}
 \Law(Z)(A)=\P(Z\in A),
 \qquad \text{for every allowed set }A\text{ of values of }Z.
 \label{eq:law}
\end{equation}
We may name this law
\[
 P_Z=\Law(Z).
\]
Thus \(\P\) acts on events of the experiment, while \(P_Z\) acts on sets in
the value space of \(Z\).

Capital letters do not determine an object's type by themselves. For example,
Lecture~1 uses \(P_{\max}\) for a position-embedding matrix and
\(Q,K,V\) for attention matrices. Their dimensions declare their types. In a
single scope, do not use the same symbol for both a probability law and a
matrix.

\section{The ambient probability measure}

The symbol \(\P\) denotes the probability measure of the stated experiment.
All random quantities in an argument are jointly defined on that experiment
unless another probability measure is explicitly introduced.

Formally, \(\P\) takes an event. The course accepts the standard inverse-image
shorthand
\begin{align}
 \P(X=x)
 &:=\P\bigl(\{\omega:X(\omega)=x\}\bigr),
 \label{eq:event-equality}\\
 \P(X\in A)
 &:=\P\bigl(\{\omega:X(\omega)\in A\}\bigr),
 \label{eq:event-membership}\\
 \P(X\in A,Y\in B)
 &:=\P\bigl(\{X\in A\}\cap\{Y\in B\}\bigr).
 \label{eq:event-joint}
\end{align}
We also use
\begin{equation}
 \P(E,F):=\P(E\cap F)
 \label{eq:event-comma}
\end{equation}
when \(E\) and \(F\) are events. The comma in \cref{eq:event-comma}
means intersection. By contrast, the comma in \(\Law(X,Y)\) constructs the
random tuple \((X,Y)\).

Expressions such as \(\P(X)\) are meaningful only when \(X\) itself denotes
an event. Prose cannot repair an ambiguous expression. In particular, a
sentence saying that a probability is ``over the initialization'' should be
unnecessary: the joint experiment and its random quantities must already make
the meaning of \(\P\) unique.

Use the course macro \verb|\P|. Do not introduce decorated probability
operators such as \(\P_{w,b}\) or \(\P_{X\sim P}\). If the experiment changes,
state the new probability measure or joint law before using probability
notation.

\section{Laws and the law operator}

The operator \(\Law\) takes random elements and returns their distribution.
In particular,
\[
 \Law(X)
 \quad\text{and}\quad
 \Law(X,Y):=\Law((X,Y))
\]
are the marginal law of \(X\) and the joint law of \((X,Y)\). The expression
\(\Law(E)\) is inappropriate when \(E\) is an event; the relevant object is
the number \(\P(E)\), or the Bernoulli law \(\Law(\I(E))\) when that law is
actually needed.

In future course material, reserve \(\mathcal L\) for this law operator. Use
another symbol, such as \(\mathcal A\), for a learning algorithm.

If a law is named \(P\), then \(P(A)\) is the probability assigned to the
measurable set \(A\). For a finite space we use
\begin{equation}
 P(x):=P(\{x\}),
 \qquad
 P(x,y):=P(\{(x,y)\}).
 \label{eq:finite-mass}
\end{equation}
If \(c\) is a measurable function on the space carrying \(P\), we also accept
the preimage shorthand
\[
 P(c=0):=P(\{z:c(z)=0\}),
 \qquad
 P(c\in A):=P(c^{-1}(A)).
\]

When \(P\) is a named law on a canonical product space and \(X,Y\) are its
coordinate maps, expressions such as
\[
 P(X=x),
 \qquad
 P(Y=y\mid X=x)
\]
are accepted shorthand for event probabilities computed under \(P\). This
notation returns numbers. It does not name the law of \((X,Y)\).

Consequently,
\[
 P_{\rm tr}(X,Y)\ne P_{\rm dep}(X,Y)
\]
should not be used to say that two joint laws differ. If the joint laws have
already been named, write
\[
 P_{\rm tr}\ne P_{\rm dep}.
\]
If they have not been named, write the law of the relevant tuple, for example
\(\Law(X,Y)\).

The established generic notation for a marginal is \(P_Z=\Law(Z)\). When a
joint law already carries a tag, introduce its marginal explicitly, for
example
\[
 P_{{\rm tr},X}
 \quad\text{or}\quad
 \mu_{\rm tr}.
\]
Do not introduce a new convention such as \(P_{\rm tr}^X\) without a reason.

\section{Mass functions, probability vectors, and predictive kernels}

On a finite set \(\mathcal X\), a probability mass function is a function
\[
 p:\mathcal X\to[0,1],
 \qquad
 \sum_{x\in\mathcal X}p(x)=1.
\]
It determines a law \(P\) by
\[
 P(A)=\sum_{x\in A}p(x).
\]
Finite-space writing often identifies the law and its mass function when the
meaning is clear. Lowercase notation emphasizes that \(p(x)\) is a number.

A probability vector is the same information with a fixed enumeration of the
outcomes. Lecture~1 writes
\[
 p=(p_v)_{v=1}^V\in\Delta_V.
\]
Here \(p_v\) is a coordinate of a deterministic vector, not an application of
the ambient probability measure \(\P\).

A predictive kernel assigns a probability distribution to every allowed
input. Lecture~1 defines
\[
 p_\theta:[V]^*\to\Delta_V,
 \qquad
 x_{<t}\longmapsto p_\theta(\,\cdot\mid x_{<t}).
\]
Therefore
\[
 p_\theta(v\mid x_{<t})
\]
is the mass assigned to token \(v\) by the distribution selected at prefix
\(x_{<t}\). The vertical bar separates the output value from the kernel input.
It does not assert that an ambient joint experiment was defined first and then
conditioned.

Sampling successively from the kernel constructs a stochastic process. The
statement
\[
 X_t\sim p_\theta(\,\cdot\mid X_{<t})
\]
is a sampling instruction. More explicitly, it requires
\[
 \Law(X_t\mid X_{<t}=x_{<t})
 =p_\theta(\,\cdot\mid x_{<t})
\]
at every prefix to which the construction assigns positive probability.
The resulting joint mass of a fixed prefix is
\[
 p_\theta(x_{1:T})
 =\prod_{t=1}^T p_\theta(x_t\mid x_{<t}).
\]

Thus \(p_\theta\) can denote a model-specified kernel or the joint mass it
induces, with the argument making the role clear. It does not replace \(\P\)
for events in a training or sampling experiment.

\section{Conditioning}

For events \(E,C\) with \(\P(C)>0\),
\begin{equation}
 \P(E\mid C)=\frac{\P(E\cap C)}{\P(C)}.
 \label{eq:conditional-probability}
\end{equation}
The notation
\[
 \P(Y=y\mid X=x)
\]
conditions the event \(\{Y=y\}\) on \(\{X=x\}\).

For jointly distributed random elements \(X,Y\),
\[
 \Law(Y\mid X=x)
\]
is a probability law on the value space of \(Y\). In the finite case, for
every set \(B\),
\[
 \Law(Y\mid X=x)(B)=\P(Y\in B\mid X=x).
\]
The notes work mainly with finite or Euclidean spaces, where the required
conditional laws are available. At values \(x\) with \(\P(X=x)=0\), the joint
law does not determine the elementary ratio in
\cref{eq:conditional-probability}. A conditional law there requires a chosen
version or a conditional target specified separately for every input.

Suppose that \(P\) and \(Q\) have been defined as kernels from the value space
of \(X\) to probability laws on the value space of \(Y\). In a discrete space,
\[
 P(Y\mid X)
 \]
is then a random variable: at outcome \(\omega\), its value is
\(P(Y(\omega)\mid X(\omega))\). Consequently,
\[
 P(Y\mid X)=Q(Y\mid X)
\]
is a valid equality of random variables, understood up to almost-sure
equality under the ambient probability measure \(\P\). It says that the two
kernels agree when evaluated at the random pair \((X,Y)\). It does not imply
that the kernels agree at values that the pair does not visit.

Similarly, for one fixed pair \((x,y)\), the scalar equality
\[
 P(y\mid x)=Q(y\mid x)
\]
compares only those two numbers. It does not assert equality elsewhere.

When the intended statement is equality of the kernels, declare their common
domain and codomain and write the preferred form
\[
 P=Q.
\]
In a discrete space, this equality expands to
\[
 P(y\mid x)=Q(y\mid x)
 \qquad\text{for every }x,y.
\]
On a general output space, it expands to
\[
 P(B\mid x)=Q(B\mid x)
 \qquad\text{for every input }x
 \text{ and every allowed output set }B.
\]

Condition only when the conditional object is used. If a calculation divides
by \(\P(C)\) and later multiplies by \(\P(C)\), a direct unconditional
factorization is usually clearer and also covers \(\P(C)=0\).

\section{Expectation}

If \(U\) is integrable, then
\[
 \E[U]=\int_\Omega U(\omega)\,\P(d\omega).
\]
Use plain \(\E[U]\) when \(U\) belongs to the ambient experiment. The setup
must identify the random quantities and their joint law before the expectation
appears.

Suppose that \(P\) is a probability law on \((\mathcal Z,\mathcal H)\). The
subscript in \(\E_P\) says to use \(P\) as the probability measure. Formally,
we may replace the ambient probability space by the canonical probability
space
\[
 (\mathcal Z,\mathcal H,P)
\]
and regard the identity map \(z\mapsto z\) as its canonical random element.
For an integrable function \(f:\mathcal Z\to\R\), the resulting expectation
is
\begin{equation}
 \E_P[f]:=\int_{\mathcal Z} f(z)\,P(dz).
 \label{eq:law-expectation}
\end{equation}
On a finite space,
\[
 \E_P[f]=\sum_z P(z)f(z).
\]
If \(Z\) is a random element in the original ambient experiment and
\(\Law(Z)=P\), then
\[
 \E_P[f]=\E[f(Z)].
\]
The notation \(\E_P[f]\) is useful when \(P\) is itself one of the objects
being compared, as in \(\E_P[c]-\E_Q[c]\).

Do not use a sampling declaration as an expectation subscript, as in
\[
 \E_{Z\sim P}[h(Z)]
 .
\]
Write \(\E_P[h]\) when the integrand is a function on the space of \(P\), or
declare \(\Law(Z)=P\) and write \(\E[h(Z)]\).

The notation \(P(dz)\) in an integral is measure-integration notation; \(dz\)
is not an event supplied as an ordinary argument to \(P\). Use this notation
only when the measure-theoretic form serves the argument. Prefer masses and
sums for finite spaces.

Conditional expectation follows the same ambient experiment:
\[
 \E[U\mid C]
 =\frac{\E[U\I(C)]}{\P(C)}
\]
when \(\P(C)>0\). The expression \(\E[U\mid X]\) is a random variable.

\section{Random and fixed quantities}

Use uppercase letters for random quantities and lowercase letters for fixed
values when this convention is available. Thus \(X=x\) denotes an event;
it does not turn \(X\) into a fixed quantity. A phrase such as ``fix \(x\)''
must precede calculations that treat \(x\) as fixed.

State joint laws and independence before using them. Separate notation does
not imply independence, and conditioning can destroy independence. When an
independence claim is used, show the corresponding factorization at the point
of use.

\section{Accepted shorthand and rejected expressions}

\begin{center}
\renewcommand{\arraystretch}{1.3}
\begin{tabular}{p{0.29\textwidth}p{0.28\textwidth}p{0.34\textwidth}}
\toprule
Expression & Status & Meaning or repair\\
\midrule
\(\P(X=x)\) & accepted & probability of the event \(\{X=x\}\)\\
\(\P(E,F)\) & accepted & \(\P(E\cap F)\) for events \(E,F\)\\
\(\P(X)\) & rejected in general & meaningful only if \(X\) denotes an event\\
\(\Law(X,Y)\) & accepted & joint law of the tuple \((X,Y)\)\\
\(\Law(E)\) & rejected for an event & use \(\P(E)\) or \(\Law(\I(E))\)\\
\(P(x)\) & accepted on a finite space & \(P(\{x\})\)\\
\(P(c=0)\) & accepted & \(P(\{z:c(z)=0\})\)\\
\(P(X=x)\) & accepted with setup & event probability under the named law \(P\)\\
\(P(X,Y)\) as a joint law & rejected & use \(P\) or \(\Law(X,Y)\)\\
\(p_\theta(v\mid x)\) & accepted & value of a model-specified predictive kernel\\
\(P(dx)\) & restricted & integration notation, introduced only when useful\\
\(\E_P[h]\) & accepted & \(\int h\,dP\)\\
\(\E_{X\sim P}[h(X)]\) & rejected by course convention & use \(\E_P[h]\), or specify \(\Law(X)=P\) and use \(\E[h(X)]\)\\
\bottomrule
\end{tabular}
\end{center}

Every shorthand must expand to a well-typed formal expression. Shorthand does
not license a probability measure to accept a random variable, or a law
operator to accept an event.

\section{Example: covariate shift}

Let \(\mathcal X\) and \(\mathcal Y\) be finite. Specify a conditional target
\[
 p^\star(\,\cdot\mid x)\in\Delta_{\mathcal Y}
 \qquad\text{for every }x\in\mathcal X,
\]
and two input laws \(\mu_{\rm tr}\) and \(\mu_{\rm dep}\). Define the training
and deployment joint laws through their masses:
\begin{equation}
 P_{\rm tr}(x,y)=\mu_{\rm tr}(x)p^\star(y\mid x),
 \qquad
 P_{\rm dep}(x,y)=\mu_{\rm dep}(x)p^\star(y\mid x).
 \label{eq:covariate-shift}
\end{equation}
There is \emph{covariate shift} when
\[
 \mu_{\rm tr}\ne\mu_{\rm dep}
\]
while the same conditional target \(p^\star\) appears in both factorizations.
The values of \(p^\star(\,\cdot\mid x)\) at inputs absent from training are
part of the specified prediction problem; they are not identified by the
training joint law.

There is \emph{general distribution shift} whenever
\[
 P_{\rm tr}\ne P_{\rm dep}.
\]
General distribution shift may change the input law, the conditional target,
or both.

\section{Checklist for new writing}

Before finalizing a probability argument, check the following.

\begin{enumerate}[leftmargin=*]
\item Every symbol has a declared type: random element, event, law, mass
function, density, probability vector, kernel, matrix, or number.
\item Every use of \(\P\) has a unique probability space and joint experiment.
\item Every displayed argument of \(\P\) expands to an event.
\item Every displayed argument of a named law expands to a measurable set, or
uses one of the finite-space or preimage shorthands above.
\item \(\Law\) is used for the law of a random element or tuple, not for an
event probability.
\item Conditional laws at zero-probability inputs are either unnecessary or
specified through an explicit version or target kernel.
\item Expectations use \(\E[U]\) for the ambient experiment and
\(\E_P[f]\) when the named law \(P\) is part of the claim.
\item Independence and changes of measure are stated before they are used.
\item Finite-space arguments use masses and sums unless measure notation adds
substantive value.
\item Every accepted shorthand can be expanded into a well-typed formal
expression.
\end{enumerate}

\end{document}
